\begin{proof}
We prove the theorem by translating the geometric hypotheses into the language of local ramification theory, using only the Cartesian property of the master diagram and the formal properties of the Herbrand function. We completely bypass Artin-Schreier-Witt theory.

\paragraph{Step 1: Linear disjointness from the Cartesian diagram.}
Let $T = \cP^{[d]_G}$, $B = (\cP_{G/H})^{[d]}$, and $B' = (\cP_{G/H})^{(d)}$. 
Consider the function fields of these schemes over the base field $k$:
\begin{itemize}
    \item $K(T)$ is the function field of $T$, with Galois group $K_G \rtimes S_d$ over $K(\cP^{\times d})$, where $K_G = \ker(G^d \to G)$.
    \item $K(B')$ is the function field of $B'$, with Galois group $H^d \rtimes S_d$.
    \item $K(B)$ is the function field of $B$, with Galois group $(K_G + H^d) \rtimes S_d$.
\end{itemize}
Because $G$ is abelian, the subgroup of $G^d \rtimes S_d$ generated by $K_G \rtimes S_d$ and $H^d \rtimes S_d$ is exactly $(K_G + H^d) \rtimes S_d$. By Galois theory, this implies that
\[
K(T) \cap K(B') = K(B).
\]
Thus, $K(T)$ and $K(B')$ are linearly disjoint over $K(B)$, and their compositum $K(T)\cdot K(B')$ is the function field of the fiber product $\cP^{[d]_H}$. This yields a canonical isomorphism of Galois groups:
\[
\operatorname{Gal}(K(\cP^{[d]_H})/K(B')) \;\xrightarrow{\;\sim\;}\; \operatorname{Gal}(K(T)/K(B)) \;=\; H.
\]

\paragraph{Step 2: The inertia gap.}
Hypothesis (H2) states that $\cP^{[d]_H} \to (\cP_{G/H})^{(d)}$ is unramified at $\eta_{C'}$. In terms of valuations, the inertia group $I(K(\cP^{[d]_H})/K(B'), v_{C'})$ is trivial. Under the isomorphism above, this maps injectively into the inertia group $I(K(T)/K(B), v_M)$, where $v_M$ is the restriction of $v_{C'}$ to $K(B)$. 

However, linear disjointness alone does not immediately force $I(K(T)/K(B), v_M)$ to be trivial. If the extension $K(B')/K(B)$ is ramified at $v_M$, the inertia group of the compositum could theoretically be a ``diagonal'' subgroup of $H \times K_{G/H}$---trivial when projected to $K_{G/H}$ (by H1) but potentially nontrivial when projected to $H$. To close this gap, we pass to the local completion and analyze the upper ramification filtration.

\paragraph{Step 3: Translation to local ramification via the inductive hypothesis.}
The broader context of this theorem is to prove the equivalence: ``$\cP^{[d]_G} \to U^{(d)}$ is unramified at $\eta_C$ $\iff$ $\cP \to U$ has ramification bounded by $d$ at $p_\infty$.'' We are currently executing the inductive step on the order of the cyclic $p$-group $G$. Thus, we may assume the equivalence holds for the quotient $G/H$ and for the subgroup $H$. 

Let $F/\widehat{K}_{p_\infty}$ be the local extension of complete discrete valuation fields corresponding to the $G$-torsor $\cP \to U$ at $p_\infty$. Since $G \cong \mathbb{Z}/p^r\mathbb{Z}$ and $H \cong \mathbb{Z}/p^s\mathbb{Z}$, the quotient is $G/H \cong \mathbb{Z}/p^{r-s}\mathbb{Z}$. By the inductive hypothesis, the geometric conditions (H1) and (H2) translate directly into bounds on the upper ramification jumps (in the classical upper numbering) of the local extensions:
\begin{itemize}
    \item \textbf{(H1)} implies that the highest upper ramification jump of the subextension $F^H / \widehat{K}_{p_\infty}$ (with Galois group $G/H$) is strictly less than $d$. Let the highest lower jump be $m_{r-s}$. Then its upper jump is $\varphi_{G/H}(m_{r-s}) < d$.
    \item \textbf{(H2)} implies that the relative extension $F / F^H$ (with Galois group $H$) has its highest upper ramification jump strictly less than $d$. Let the highest lower jump of the full extension $F/\widehat{K}_{p_\infty}$ be $m_r$. The corresponding upper jump for $H$ is $\varphi_H(m_r) < d$.
\end{itemize}
Our goal is to show that the highest upper ramification jump of the full extension $F/\widehat{K}_{p_\infty}$ for the group $G$ is strictly less than $d$, i.e., $\varphi_G(m_r) < d$. By the inductive hypothesis (applied in reverse), this will imply that $\cP^{[d]_G} \to U^{(d)}$ is unramified at $\eta_C$.

\paragraph{Step 4: Convexity of the Herbrand function.}
Since $G$ is cyclic, the lower ramification groups form a chain through every subgroup of $G$, with jumps at integers $m_1 < m_2 < \dots < m_r$. Let $\Delta_j = m_{j+1} - m_j \geq 1$. A direct computation of the Herbrand functions $\varphi(u) = \int_0^u \frac{|G_t|}{|G_0|}\,dt$ yields:
\[
\varphi_{G/H}(m_{r-s}) = \sum_{j=0}^{r-s-1} p^{-j}\Delta_j,
\]
\[
\varphi_H(m_r) = m_{r-s} + \sum_{k=0}^{s-1} p^{-k}\Delta_{r-s+k},
\]
\[
\varphi_G(m_r) = \sum_{j=0}^{r-1} p^{-j}\Delta_j.
\]
Splitting the sum for $\varphi_G(m_r)$ at $j = r-s$ and factoring out $p^{-(r-s)}$ from the second part gives the fundamental identity:
\[
\varphi_G(m_r) \;=\; \varphi_{G/H}(m_{r-s}) \;+\; \frac{\varphi_H(m_r) - m_{r-s}}{p^{r-s}}.
\]
Let $A = \varphi_{G/H}(m_{r-s})$ and $B = \varphi_H(m_r)$. By our translated hypotheses, $A < d$ and $B < d$. Furthermore, since $\varphi(u) \leq u$ for all $u \geq 0$, we have $A \leq m_{r-s}$. Substituting this into the identity:
\[
\varphi_G(m_r) \;=\; A + \frac{B - m_{r-s}}{p^{r-s}} \;\leq\; A + \frac{B - A}{p^{r-s}} \;=\; A\left(1 - \frac{1}{p^{r-s}}\right) + \frac{B}{p^{r-s}}.
\]
Since both $A < d$ and $B < d$, the right-hand side is a strict convex combination of two quantities strictly less than $d$:
\[
\varphi_G(m_r) \;<\; d\left(1 - \frac{1}{p^{r-s}}\right) + \frac{d}{p^{r-s}} \;=\; d.
\]
Because the upper ramification group $G^u$ is trivial for all $u > \varphi_G(m_r)$, and we have shown $\varphi_G(m_r) < d$, it follows that $G^d = \{1\}$. This means the original $G$-torsor $\cP \to U$ has ramification bounded by $d$ at $p_\infty$, which by the global-to-local equivalence for symmetric powers implies that $\cP^{[d]_G} \to U^{(d)}$ is unramified at $\eta_C$.
\end{proof}